\section{以\texorpdfstring{$a$}{a}为底的指数函数和对数函数的性态}

\begin{proposition}{}{}
	设$a\in\R_{>0}\setminus\left\{1\right\}$,$f_{a}$是以$a$为底的指数函数.
	\begin{enumerate}
		\item 若$a\in\left(1,+\infty\right)$,则$f_{a}$和$\log_{a}$严格递增,并且
		\begin{align*}
			\limit_{x\to+\infty}a^{x}=+\infty,\limit_{x\to-\infty}a^{x}=0,\limit_{\substack{x\to 0\\x\in\R_{>0}}}\log_{a}x=-\infty,\limit_{x\to+\infty}\log_{a}x=+\infty.
		\end{align*}
		\item 若$a\in\left(0,1\right)$,则$f_{a}$和$\log_{a}$严格递减,并且
		\begin{align*}
			\limit_{x\to+\infty}a^{x}=0,\limit_{x\to-\infty}a^{x}=+\infty,\limit_{\substack{x\to 0\\x\in\R_{>0}}}\log_{a}x=+\infty,\limit_{x\to+\infty}\log_{a}x=-\infty.
		\end{align*}
	\end{enumerate}
\end{proposition}

\begin{corollary}{}{}
	设$a,b\in\R_{>0}$且$a<b$.
	\begin{enumerate}
		\item 对每个$x\in\R_{>0}$,有$a^{x}<b^{x}$.
		\item 对每个$x\in\R_{<0}$,有$a^{x}>b^{x}$.
	\end{enumerate}
\end{corollary}

\begin{proposition}{实指数函数的连续延拓}{}
	设$a\in\R,f:\R\to\R_{>0},x\mapsto a^{x}$.
	\begin{enumerate}
		\item 当$a\in\left(1,+\infty\right)$时,$\tilde{f}:\left[-\infty,+\infty\right]\to\left[-\infty,+\infty\right],x\mapsto
		\begin{cases}
			-\infty,&x=-\infty\\
			f\left(x\right),&-\infty<x<+\infty\\
			+\infty,&x=+\infty
		\end{cases}$是$f$的连续延拓.
		\item 当$a\in\left(0,1\right)$时,$\tilde{f}:\left[-\infty,+\infty\right]\to\left[-\infty,+\infty\right],x\mapsto
		\begin{cases}
			+\infty,&x=-\infty\\
			f\left(x\right),&-\infty<x<+\infty\\
			-\infty,&x=+\infty
		\end{cases}$是$f$的连续延拓.
	\end{enumerate}
\end{proposition}

\begin{proposition}{实指数函数的连续延拓}{}
	设$a\in\R_{>0},f:\R_{>0}\to\R,x\mapsto\log_{a}\left(x\right)$.
	\begin{enumerate}
		\item 当$a\in\left(1,+\infty\right)$时,$\tilde{f}:\left[0,+\infty\right]\to\left[-\infty,+\infty\right],x\mapsto
		\begin{cases}
			-\infty,&x=0\\
			f\left(x\right),&-\infty<x<+\infty\\
			+\infty,&x=+\infty
		\end{cases}$是$f$的连续延拓.
		\item 当$a\in\left(0,1\right)$时,$\tilde{f}:\left[0,+\infty\right]\to\left[-\infty,+\infty\right],x\mapsto
		\begin{cases}
			+\infty,&x=0\\
			f\left(x\right),&-\infty<x<+\infty\\
			0,&x=+\infty
		\end{cases}$是$f$的连续延拓.
	\end{enumerate}
\end{proposition}

\begin{proposition}{实指数幂函数的连续延拓}{}
	设$f:\R_{>0}\times\R\to\R,\left(x,y\right)\mapsto x^{y}$.定义$\tilde{f}:\left[0,+\infty\right]\times\left[-\infty,+\infty\right]$如下:
	\begin{align*}
		\tilde{f}\left(x,y\right)=
		\begin{cases}
			f\left(x,y\right),&x\in\R_{>0},y\in\R\\
			0,&x=0,y\in\left(0,+\infty\right]\\
			+\infty,&x=0,y\in\left[-\infty,0\right)\\
			+\infty,&x=+\infty,y\in\left(0,+\infty\right]\\
			0,&x=+\infty,y\in\left[-\infty,0\right)\\
			0,&x\in\left[0,1\right),y=+\infty\\
			+\infty,&x\in\left(1,+\infty\right],y=+\infty\\
			+\infty,&x\in\left(0,1\right],y=-\infty\\
			0,&x\in\left(1,+\infty\right],y=-\infty
		\end{cases}
	\end{align*}
	那么$\tilde{f}$是$f$的连续延拓.
\end{proposition}

\begin{remark}{}{}
	\begin{enumerate}
		\item 只要等式两边都有意义,那么如下公式对于连续延拓后的实指数函数,实对数函数和实指数幂函数依然成立:
		\begin{align*}
			a^{x+y}=a^{x}a^{y},\log_{a}\left(xy\right)=\log_{a}\left(x\right)+\log_{a}\left(y\right),\left(a^{x}\right)^{y}=a^{xy}.
		\end{align*}
		\item $f:\R_{>0}\times\R\to\R,\left(x,y\right)\mapsto x^{y}$不能给出代数中约定的$0^0=1$.
		\item 当$a\in\R_{>0}$时,以$a$为底的实指数函数把$\Z\to\R,n\mapsto a^{n}$延拓到了$\R$,但是当$a\in\R_{<0}$时我们得不到任何延拓——这种函数的``自然''延拓只能在解析函数理论中定义。
	\end{enumerate}
\end{remark}